\phantomsection
\addcontentsline{toc}{subsection}{Unidad 2. Oxidación-Reducción}
\subsection*{\textcolor{blueunam}{Unidad 2. Oxidación-Reducción}}

\phantomsection
\addcontentsline{toc}{subsubsection}{Clase 10 -- 29/09/2026: Estados de oxidación en compuestos orgánicos}
\subsubsection*{\textcolor{blueunam}{Clase 10 -- 29/09/2026: Estados de oxidación en compuestos orgánicos}}

\textbf{Número de oxidación (n.o.):} conteo formal que finge que los electrones 
del enlace se los lleva \textbf{completamente} el átomo más electronegativo 
(reparto 100/0). No es una carga real ni un promedio: es una asignación contable 
para rastrear electrones en reacciones redox.

\textbf{Cómo se calcula el n.o. del carbono:} se cuenta la contribución de cada 
enlace del C.

\begin{center}
\begin{tabular}{|l|c|}
\hline
\textbf{Enlace del C} & \textbf{Contribución al n.o.} \\
\hline
A un átomo más electronegativo (O, N, halógeno), enlace simple & $+1$ \\
A un átomo más electronegativo, enlace doble (C=O) & $+2$ \\
A un átomo más electronegativo, enlace triple (C$\equiv$N) & $+3$ \\
A un átomo menos electronegativo (H) & $-1$ \\
A otro C (simple, doble o triple) & $0$ \\
\hline
\end{tabular}
\end{center}

La suma de contribuciones da el n.o. del carbono. En química orgánica estándar 
va de $-4$ a $+4$, con todos los enteros intermedios.

\begin{center}
\begin{tabular}{lccc}
\hline
\textbf{Compuesto} & \textbf{Enlaces del C} & \textbf{n.o. del C} \\
\hline
\ce{CH4}          & 4 C--H                        & $-4$ \\
\ce{CH3-CH3}      & 3 C--H, 1 C--C                & $-3$ \\
\ce{CH3OH}        & 3 C--H, 1 C--O                & $-2$ \\
\ce{CH3-CH2OH}    & 2 C--H, 1 C--C, 1 C--O        & $-1$ \\
\ce{CH2O}         & 2 C--H, 1 C=O                 & $0$  \\
\ce{CH3-CHO}      & 1 C--H, 1 C--C, 1 C=O         & $+1$ \\
\ce{HCOOH}        & 1 C--H, 1 C=O, 1 C--O         & $+2$ \\
\ce{CH3-COOH}     & 1 C--C, 1 C=O, 1 C--O         & $+3$ \\
\ce{CO2}          & 2 C=O                         & $+4$ \\
\hline
\end{tabular}
\end{center}
(Esta tabla y cómo calcular sus n.o. se desarrolla mejor más adelante.) \\

{\small\itshape Nota: en \ce{CH3-CH2OH}, \ce{CH3-CHO} y \ce{CH3-COOH}, el n.o. 
indicado corresponde al carbono señalado (el que tiene el grupo funcional), 
no al \ce{CH3} terminal.}

\textbf{Oxidación:} pérdida de \ce{e^{\color{anblue}-}}. En química orgánica, 
se reduce a ganancia de oxígeno \ce{O} o halógenos, o pérdida de hidrógenos \ce{H}. \\

\textbf{Reducción:} ganancia de \ce{e^{\color{anblue}-}}. En química orgánica, 
se reduce a pérdida de oxígeno \ce{O} o halógenos, o ganancia de hidrógenos \ce{H}. \\

Generalmente, si el \ce{C} se enlaza con átomos más electronegativos, hablamos de 
oxidación. Si se enlaza con átomos menos electronegativos, hablamos de reducción. 
(Tomando como referencia el \ce{C}). \\

\textbf{Ejemplos:}
\begin{center}
\begin{tabular}{l@{\qquad}l}
$\ce{Fe^{\color{catred}{2+}}} \longrightarrow \ce{Fe^{\color{catred}{3+}}} + \ce{e^{\color{anblue}-}}$
& (oxidación: perdió 1 \ce{e^{\color{anblue}-}}, carga $+2 \to +3$) \\[8pt]
$\ce{Cu^{\color{catred}{2+}}} + 2\,\ce{e^{\color{anblue}-}} \longrightarrow \ce{Cu}$
& (reducción: ganó 2 \ce{e^{\color{anblue}-}}, carga $+2 \to 0$) \\[8pt]
$\ce{Zn} \longrightarrow \ce{Zn^{\color{catred}{2+}}} + 2\,\ce{e^{\color{anblue}-}}$
& (oxidación: perdió 2 \ce{e^{\color{anblue}-}}, carga $0 \to +2$) \\[8pt]
$\ce{Cl2} + 2\,\ce{e^{\color{anblue}-}} \longrightarrow 2\,\ce{Cl^{\color{anblue}-}}$
& (reducción: ganó 1 \ce{e^{\color{anblue}-}} por átomo, carga $0 \to -1$)
\end{tabular}
\end{center}

\textbf{Aclaración sobre el signo de la carga:} la carga no indica directamente el 
número de electrones, sino el \textbf{balance} entre protones y electrones 
\textbf{respecto al átomo neutro}.

\begin{center}
\begin{tabular}{lcc}
\hline
\textbf{Especie} & \textbf{Carga} & \textbf{Significado} \\
\hline
\ce{Fe^{\color{catred}{3+}}} & $+3$ & 3 protones más que electrones (perdió 3 \ce{e^{\color{anblue}-}} respecto al neutro) \\
\ce{Fe^{\color{catred}{2+}}} & $+2$ & 2 protones más que electrones (perdió 2 \ce{e^{\color{anblue}-}}) \\
\ce{Fe}      & $0$  & Igual número de protones y electrones \\
\ce{Cl^{\color{anblue}-}} & $-1$ & 1 electrón más que protones (ganó 1 \ce{e^{\color{anblue}-}}) \\
\hline
\end{tabular}
\end{center}

Un número más positivo significa \textbf{menos electrones}. Un número más negativo 
significa \textbf{más electrones}. Esto no es intuitivo a primera vista, pero hay 
que recordar que los electrones tienen carga negativa: si el átomo pierde 
electrones (que son negativos), su carga se vuelve más positiva. Por eso $+3$ no 
significa ``3 electrones extra'', sino ``3 electrones de menos''.

\begin{center}
\begin{tikzpicture}[xscale=1.3]
    \draw[-{Stealth}, thick] (-4.8,0) -- (4.8,0);

    \foreach \x in {-4,...,4}
        \draw[thick] (\x,0.18) -- (\x,-0.18);

    \foreach \x in {-4,...,4}
        \node[font=\small, fill=white, inner sep=1.5pt] at (\x,0) {\x};

    \node[left] at (-5.1,0) {\small \textbf{red}};
    \node[right] at (5.1,0) {\small \textbf{oxi}};
\end{tikzpicture}
\end{center}

\textbf{Ejemplos y su posición en la escala:} \\

\textbf{Pasos:}
\begin{enumerate}
    \item Identificar el carbono de interés.
    \item Identificar los sustituyentes directos a su alrededor.
    \item Átomos más electronegativos que el \ce{C}: se les asigna carga 
    $-1$ por cada enlace al carbono, y al carbono se le suma $+1$ por cada 
    enlace con estos.
    \item Átomos menos electronegativos que el \ce{C}: se les asigna carga 
    $+1$ por cada enlace al carbono, y al carbono se le suma $-1$ por cada 
    enlace con estos.
    \item Enlaces \ce{C-C}: aportan $0$ a ambos carbonos.
\end{enumerate} 

\textbf{Nota:} para calcular el n.o. del carbono de interés solo cuentan 
sus enlaces directos. Lo que ocurra con los átomos más allá del primero 
(el \ce{H} de un \ce{O-H}, por ejemplo) no afecta al carbono. Esos enlaces 
solo se consideran cuando se balancea la carga total de la molécula, 
sumando los n.o. de todos los átomos. \\

\textbf{Nota:} en moléculas con varios carbonos con n.o. distinto, 
la escala registra el \textbf{carbono de interés} (el del grupo funcional), 
no un promedio ni la molécula entera. Los demás carbonos tienen su propio 
n.o. que se calcula por separado. \\

\ce{CH4} (Grupo: Alcano) -- Metano: \\
\begin{center}
\schemestart
\chemfig{@{C}C(-[:0]H)(-[:90]H)(-[:180]H)(-[:270]H)}
\schemestop
\chemmove{
    \draw[-{Stealth}, thick] (C) -- +(2,1.5)
        node[right, text width=5cm]{\small El \ce{C} es más electronegativo que el \ce{H}: se lleva los 4\,\ce{e^{\color{anblue}-}}, uno de cada enlace \ce{C-H}};
}
\end{center}
Por lo que para este, el \ce{C} tiene un n.o. de $-4$: \ce{C^{\color{anblue}{4-}}} y cada
\ce{H} tiene n.o. de $+1$: \ce{H^{\color{catred}+}}. En la escala se
posiciona en el estado más reducido: n.o. $= -4$.

\ce{CH3Cl} (Grupo: Haloalcano) -- Clorometano: \\
\begin{center}
\schemestart
\chemfig{C^{\color{anblue}{2-}}(-[:0]Cl^{\color{anblue}-})(-[:90]H^{\color{catred}+})(-[:180]H^{\color{catred}+})(-[:270]H^{\color{catred}+})}
\schemestop
\end{center}
Verificación: $(-2) + 3(+1) + (-1) = 0$. En la escala se posiciona en $-2$.

\ce{CO2} (Grupo: Óxido inorgánico) -- Dióxido de carbono: \\
\begin{center}
\schemestart
\chemfig{C^{\color{catred}{4+}}(=[:0]O^{\color{anblue}{2-}})(=[:180]O^{\color{anblue}{2-}})}
\schemestop
\end{center}
Cada enlace del \ce{C} con un átomo más electronegativo le aporta $+1$ al 
carbono y $-1$ al otro átomo. En el \ce{CO2} cada \ce{O} tiene un enlace 
doble con el \ce{C}, es decir dos enlaces: $2(-1) = -2$, y el \ce{C} queda 
en $+4$. Verificación: $(+4) + 2(-2) = 0$. En la escala se posiciona en 
el estado más oxidado: n.o. $= +4$.

\ce{HCHO / CH2O} (Grupo: Aldehído) -- Metanal / Formaldehído: \\
\begin{center}
\schemestart
\chemfig{C^0(=[:90]O^{\color{anblue}2-})(-[:230]H^{\color{catred}+})(-[:310]H^{\color{catred}+})}
\schemestop
\end{center}
Verificación: $(0) + 2(+1) + (-2) = 0$. En la escala se posiciona en $0$.

\ce{CH2=CH2 / C2H4} (Grupo: Alqueno) -- Eteno / Etileno: \\
\begin{center}
\schemestart
\chemfig{C^{\color{anblue}{2-}}(-[:140]H^{\color{catred}+})(-[:230]H^{\color{catred}+})=C^{\color{anblue}{2-}}(-[:40]H^{\color{catred}+})(-[:320]H^{\color{catred}+})}
\schemestop
\end{center}
Verificación: $2(-2) + 4(+1) = 0$. En la escala se posiciona en $-2$.

\ce{CH3OH / CH4O} (Grupo: Alcohol) -- Metanol / Alcohol metílico: \\
\begin{center}
\schemestart
\chemfig{C^{\color{anblue}{2-}}(-[:0]O^{\color{anblue}{2-}}H^{\color{catred}+})(-[:90]H^{\color{catred}+})(-[:180]H^{\color{catred}+})(-[:270]H^{\color{catred}+})}
\schemestop
\end{center}
Verificación: $(-2) + 4(+1) + (-2) = 0$. En la escala se posiciona en $-2$.

\ce{HCOOH / CH2O2} (Grupo: Ácido carboxílico) -- Ácido metanoico / Ácido fórmico: \\
\begin{center}
\schemestart
\chemfig{C^{\color{catred}{2+}}(=[:90]O^{\color{anblue}2-})(-[:180]H^{\color{catred}+})(-[:0]O^{\color{anblue}2-}(-[:0]H^{\color{catred}+}))}
\schemestop
\end{center}
Verificación: $(+2) + 2(+1) + 2(-2) = 0$. En la escala se posiciona en $+2$.

\ce{HCN} (Grupo: Nitrilo) -- Formonitrilo: \\
\begin{center}
\schemestart
\chemfig{C^{\color{catred}{2+}}(-[:180]H^{\color{catred}+})(~[:0]N^{\color{anblue}{3-}})}
\schemestop
\end{center}
Cada enlace del \ce{C} con un átomo más electronegativo le aporta $+1$ al 
carbono y $-1$ al otro átomo. El \ce{N} tiene un enlace triple con el \ce{C}, 
es decir tres enlaces: $3(-1) = -3$. El \ce{H} es menos electronegativo que 
el \ce{C}, así que el enlace C--H aporta $-1$ al \ce{C} y $+1$ al \ce{H}. 
Sumando: el \ce{C} queda en $+3 - 1 = +2$. Verificación: 
$(+2) + (-3) + (+1) = 0$. En la escala se posiciona en el estado 
$+2$.

\ce{CH3CH2OCH2CH3 / C4H10O} (Grupo: Éter) -- Etoxietano / Dietil éter: \\
\begin{center}
\schemestart
\chemfig{[:-45]-[:45]-O-[:45]-[:-45]}
\qquad\qquad
\chemfig{
[:-45]C^{\color{anblue}{3-}}(-[:140]H^{\color{catred}+})(-[:220]H^{\color{catred}+})(-[:310]H^{\color{catred}+})
-[:45]C^{\color{anblue}{-}}(-[:90]H^{\color{catred}+})(-[:270]H^{\color{catred}+})
-O^{\color{anblue}{2-}}
-[:45]C^{\color{anblue}{-}}(-[:90]H^{\color{catred}+})(-[:270]H^{\color{catred}+})
-[:-45]C^{\color{anblue}{3-}}(-[:40]H^{\color{catred}+})(-[:220]H^{\color{catred}+})(-[:310]H^{\color{catred}+})
}
\schemestop
\end{center}
Verificación: $2(-3) + 2(-1) + (-2) + 10(+1) = 0$. En la escala se posiciona en $-1$ (C unido al \ce{O}). \\
\textbf{Nota:} los enlaces C--C no aportan carga (paso 5), por eso cada 
carbono puede tener n.o. distinto dentro de la misma molécula. Conviene 
distinguir entre el n.o. de un carbono individual y la suma de todos los 
n.o., que en una molécula neutra da $0$. Además existe otro concepto: la 
suma de las cargas de todos los carbonos, que indica cuánta carga aportan 
los demás átomos para balancear. Ese segundo concepto se desarrolla 
más adelante.

\ce{CCl4} (Grupo: Haloalcano) -- Tetraclorometano: \\
\begin{center}
\schemestart
\chemfig{C^{\color{catred}{4+}}(-[:0]Cl^{\color{anblue}-})(-[:90]Cl^{\color{anblue}-})(-[:180]Cl^{\color{anblue}-})(-[:270]Cl^{\color{anblue}-})}
\schemestop
\end{center}
Verificación: $(+4) + 4(-1) = 0$. En la escala se posiciona en $+4$.

\ce{CH3COCH3 / C3H6O} (Grupo: Cetona) -- Propanona / Acetona: \\
\begin{center}
\schemestart
\chemfig{C^{\color{catred}{2+}}(=[:90]O^{\color{anblue}{2-}})(-[:225]^{\color{anblue}{3-}})(-[:315]^{\color{anblue}{3-}})}
\schemestop
\end{center}
Verificación: $(+2) + 2(-3) + (-2) + 6(+1) = 0$. En la escala se posiciona en $+2$.

\ce{C6H12} (Grupo: Cicloalcano) -- Ciclohexano: \\
\begin{center}
\schemestart
\chemfig{*6(^{\color{anblue}{2-}}-^{\color{anblue}{2-}}-^{\color{anblue}{2-}}-^{\color{anblue}{2-}}-^{\color{anblue}{2-}}-^{\color{anblue}{2-}}-)}
\schemestop
\end{center}
Verificación: $6(-2) + 12(+1) = 0$. En la escala se posiciona en $-2$.

\ce{C6H6} (Grupo: Areno) -- Benceno: \\
\begin{center}
\schemestart
\chemfig{*6(^{\color{anblue}{-}}=^{\color{anblue}{-}}-^{\color{anblue}{-}}=^{\color{anblue}{-}}-^{\color{anblue}{-}}=^{\color{anblue}{-}}-)}
\schemestop
\end{center}
Verificación: $6(-1) + 6(+1) = 0$. En la escala se posiciona en $-1$.

\begin{center}
\begin{tikzpicture}[xscale=1.5, yscale=1.5]
    \draw[-{Stealth}, thick] (-4.2,0) -- (4.2,0);

    \foreach \x in {-4,...,4}
        \draw[thick] (\x,0.15) -- (\x,-0.15);

    \foreach \x in {-4,...,4}
        \node[font=\small, fill=white, inner sep=1.5pt] at (\x,0) {\x};

    \node[left] at (-4.5,0) {\small \textbf{red}};
    \node[right] at (4.5,0) {\small \textbf{oxi}};

    % -4
    \node[above=10pt] at (-4,0) {\small \ce{CH4}};

    % -2 (cuatro compuestos apilados)
    \node[above=10pt] at (-2,0) {\small \ce{CH3OH}};
    \node[above=28pt] at (-2,0) {\small \ce{CH3Cl}};
    \node[below=10pt] at (-2,0) {\small \ce{C2H4}};
    \node[below=28pt] at (-2,0) {\small \ce{C6H12}};

    % -1
    \node[above=10pt] at (-1,0) {\small \ce{C6H6}};
    \node[below=10pt] at (-1,0) {\small éter};

    % 0
    \node[above=10pt] at (0,0) {\small \ce{HCHO}};

    % +2 (tres compuestos apilados)
    \node[above=10pt] at (2,0) {\small \ce{HCOOH}};
    \node[below=10pt] at (2,0) {\small \ce{HCN}};
    \node[below=28pt] at (2,0) {\small Acetona};

    % +4
    \node[above=10pt] at (4,0) {\small \ce{CO2}};
    \node[below=10pt] at (4,0) {\small \ce{CCl4}};
\end{tikzpicture}
\end{center}

Una vez vistos varios ejemplos se pueden notar dos cosas. Imaginemos la 
reacción de combustión del metano:

\begin{center}
\schemestart
\ce{CH4}
\+
\ce{2O2}
\arrow{->[$\Delta$]}
\ce{CO2}
\+
\ce{2H2O}
\schemestop
\end{center}

Es un proceso de oxidación del \ce{C^{\color{anblue}{4-}}} al 
\ce{C^{\color{catred}{4+}}}, con un cambio neto de 8 electrones en un solo 
paso. Sin embargo, esto no significa que todas las oxidaciones sean así. 
La serie de compuestos que va del metano al dióxido de carbono muestra que 
la oxidación del carbono puede darse en saltos de 2 electrones cuando el 
proceso se controla paso a paso:

\begin{center}
\schemestart
\ce{CH4}
\arrow{->[$-2\ce{e-}$]}
\ce{CH3OH}
\arrow{->[$-2\ce{e-}$]}
\ce{CH2O}
\arrow{->[$-2\ce{e-}$]}
\ce{HCOOH}
\arrow{->[$-2\ce{e-}$]}
\ce{CO2}
\schemestop
\\[8pt]
{\small $n.o.$: \quad $-4$ \quad $\longrightarrow$ \quad $-2$ \quad 
$\longrightarrow$ \quad $0$ \quad $\longrightarrow$ \quad $+2$ \quad 
$\longrightarrow$ \quad $+4$}
\end{center}

\textbf{Para qué sirve el n.o.:}

\begin{enumerate}
    \item \textbf{Reactividad.} El n.o. indica si el C es pobre o rico en 
    densidad electrónica, y por lo tanto qué lo ataca. Un n.o. alto 
    significa que el C está deficiente de electrones y tiene carga parcial 
    positiva, así que se comporta como electrófilo y lo atacan los 
    nucleófilos. El \ce{CO2}, con su C en $+4$, reacciona con 
    \ce{OH^{\color{anblue}{-}}} para dar bicarbonato. En cambio, un n.o. 
    bajo significa que el C está saturado de electrones y no tiene nada 
    que un nucleófilo quiera atacar. El \ce{CH4}, con su C en $-4$, no 
    reacciona con \ce{OH^{\color{anblue}{-}}}, pero sí con oxidantes como 
    el \ce{O2}, que arrancan electrones de donde sobran. Un n.o. intermedio 
    tiene ambas rutas. El metanal, por ejemplo, puede oxidarse a ácido 
    fórmico o reducirse a metanol según las condiciones. El n.o. da la 
    dirección termodinámica pero no la velocidad. El \ce{CO2} podría 
    reaccionar con agua, pero la barrera de activación es alta y a 
    temperatura ambiente no lo hace.

    \item \textbf{Energía almacenada.} Un enlace fuerte cuesta mucho 
    romperlo, mientras que uno débil cuesta poco. La fuerza viene de la 
    diferencia de electronegatividad entre los átomos. Cuando un átomo muy 
    electronegativo, como el O, atrae los electrones compartidos, el enlace 
    se polariza y gana una estabilización electrostática extra. En cambio, 
    cuando los dos átomos tienen electronegatividad parecida, como C y H o 
    C y C, la compartición es casi equitativa y el enlace queda débil.

    Los enlaces alrededor de un C reducido, \ce{C-H} y \ce{C-C}, son 
    débiles y tienen energías de $\sim$413 y $\sim$347 kJ\,mol$^{-1}$. 
    Los de un C oxidado, \ce{C=O} y \ce{O-H}, son fuertes y tienen 
    $\sim$745 y $\sim$463 kJ\,mol$^{-1}$.

    Al oxidar un compuesto se rompen enlaces débiles y se forman enlaces 
    fuertes. Formar un enlace fuerte libera más energía de la que costó 
    romper uno débil, y esa diferencia sale como calor. Por eso los 
    combustibles son moléculas reducidas como el metano, el octano o la 
    glucosa, mientras que el \ce{CO2} es el producto final inerte de la 
    combustión y la respiración: ya está en el fondo del pozo energético.

    \item \textbf{Reacciones posibles.} El C tiene cuatro electrones de 
    valencia, ni más ni menos. Por eso los n.o. extremos son $+4$ y $-4$. 
    En el \ce{CO2} el carbono perdió los cuatro, y en el \ce{CH4} los 
    ganó. No existen n.o. $+5$ o $-5$ estables. Del $-4$ solo se puede 
    oxidar, porque ya no hay más electrones que ganar. Del $+4$ solo se 
    puede reducir, porque ya no hay más electrones que perder. Del 
    intermedio se pueden hacer ambas. Con esto se predice qué transformación 
    es posible antes de intentar la reacción. Un compuesto en $+3$, como 
    el ácido acético \ce{CH3-COOH}, puede oxidarse a $+4$ o reducirse 
    hacia abajo, pero no puede saltar directo a $-4$ en un solo paso 
    concertado.

    \item \textbf{Equivalencia entre compuestos.} La serie \ce{CH4} 
    ($-4$) $\to$ \ce{CH3OH} ($-2$) $\to$ \ce{CH2O} ($0$) $\to$ \ce{HCOOH} 
    ($+2$) $\to$ \ce{CO2} ($+4$) avanza de 2 en 2 electrones porque ningún 
    compuesto tiene enlaces C--C. Cuando aparecen enlaces C--C, también 
    aparecen valores impares como $-3$, $-1$, $+1$ y $+3$. El n.o. permite 
    comparar dos compuestos cualesquiera de la serie, no solo los vecinos. 
    Por ejemplo, de \ce{CH4} a \ce{HCOOH} se transfieren 6 electrones. Esto 
    sirve para balancear semirreacciones redox o para calcular cuánto 
    oxidante se necesita en una reacción.

    \textbf{Por qué los saltos son de 2 electrones.} Los enlaces covalentes 
    contienen pares de electrones. Cuando un enlace C--H se convierte en 
    C--OH o en C=O, los dos electrones del enlace original se reorganizan 
    para formar el nuevo enlace con el heteroátomo. En esta serie el 
    heteroátomo es siempre el oxígeno, pero la regla es general: también 
    aplica con N, Cl o cualquier átomo más electronegativo que el C. Cada 
    heteroátomo incorporado representa un par de electrones transferidos, 
    y la oxidación avanza de 2 en 2. En las reacciones polares comunes, 
    como la sustitución nucleofílica, la eliminación o la adición a un 
    carbonilo, las flechas curvas siempre representan el movimiento de un 
    par de electrones, nunca de uno solo.

    Los saltos de 1 electrón sí existen, pero generan radicales, especies 
    con un electrón desapareado y de vida corta. Cuando ocurre una 
    transferencia de un solo electrón, la reacción deja de ser polar y se 
    vuelve radicalaria por definición. Esto pasa en la formación de 
    reactivos de Grignard, en la reducción de Birch, en fotoredox y en la 
    combustión, donde las cadenas radicalarias propagan la reacción 
    quitando hidrógenos de a uno. También aparece con metales de transición 
    que cambian de estado de oxidación en una unidad, como el par 
    \ce{Fe^{\color{catred}{2+}}}/\ce{Fe^{\color{catred}{3+}}}, porque el hierro 
    puede ceder o aceptar electrones de uno en uno. En bioquímica, en 
    cambio, la cadena transportadora de electrones mueve los electrones de 
    a pares a través de deshidrogenasas que los transfieren al \ce{NAD+} 
    o al \ce{FAD}, manteniendo la lógica de 2 en 2.
\end{enumerate}

\textbf{Suma de n.o. de todos los carbonos.} Hasta aquí cada cálculo fue 
para un carbono individual, el del grupo funcional. Pero también se puede 
sumar el n.o. de \textbf{todos} los carbonos de la molécula. Esa suma 
indica la carga que aporta el esqueleto carbonado, y el resto de los 
átomos (H, O, N, halógenos) suman lo opuesto para balancear.

\textbf{Regla:} la suma de los n.o. de \emph{todos} los átomos es igual a 
la carga total de la especie. En una molécula neutra da $0$, en un catión 
da la carga positiva, en un anión da la carga negativa.

\textbf{Atajo:} contar los H y los heteroátomos es más rápido que analizar 
cada carbono. Lo que falte para llegar a la carga total es lo que aportan 
los carbonos.

\textbf{Ejemplo 1: isopropanol} \ce{CH3-CH(OH)-CH3}.

\begin{itemize}
    \item 7 \ce{H} unidos a C: $7(+1) = +7$
    \item 1 \ce{O} unido solo al C: $-1$
    \item El \ce{H} del \ce{O-H} no cuenta: junto con el \ce{O} suman $+1 - 1 = 0$
    \item Total H + O: $+6$
    \item Molécula neutra, así que la suma de C debe ser $-6$
\end{itemize}

Verificación por carbono: C1 (CH$_3$) $= -3$, C2 (CH) $= 0$, C3 (CH$_3$) $= -3$. Suma: $-6$.

\textbf{Ejemplo 2: catión alílico} \ce{CH2=CH-CH2+}.

\begin{itemize}
    \item 5 \ce{H} unidos a C: $5(+1) = +5$
    \item Sin heteroátomos
    \item Carga total del ion: $+1$
    \item Suma de C: $-(+5) + (+1) = -4$
\end{itemize}

Verificación por carbono: C1 (CH$_2$=) $= -2$, C2 (CH=) $= -1$, C3 (CH$_2^+$) $= -1$. Suma: $-4$.

\textbf{Utilidad en moléculas grandes.} En una molécula con muchos 
carbonos, calcular el n.o. de cada uno por separado es tedioso. La suma 
total de carbonos se obtiene de una sola cuenta con los H y heteroátomos, 
y comparar esa suma antes y después de una reacción da directamente los 
electrones transferidos. En materiales como los COFs, donde cada nodo 
repite el patrón miles de veces, este atajo es la forma práctica de 
rastrear la carga en la red sin analizar cada carbono individual.

\textbf{Ejemplo de COFs:}
\begin{center}
\schemestart
\chemfig{
    *6((-F)=(
        -N=[:225]-[:270](*6(=-=-=(
            -[:45]NH-[:-45]-[:45](*6(=-=-(
                -[:90]N=[:45]-[:90](*6(=-=-(
                    -[:135]=[:225]N?
                )=-))
            )=-))
        )-))
    )-(-F)=?-(-F)=(-F)-(-F))
}
\schemestop
\end{center}

\textbf{Análisis de n.o. en el fragmento de COF.} El fragmento tiene 
cuatro tipos de carbono distinguibles por sus sustituyentes.

\begin{center}
\small
\renewcommand{\arraystretch}{1.5}
\begin{tabular}{lccc}
\hline
\textbf{Tipo} & \textbf{Enlaces del C} & \textbf{n.o.} & \textbf{Ubicación} \\
\hline
C--F aromático    & 2 C--C, 1 C--F        & $+1$ & anillo fluorado \\
C--H aromático    & 2 C--C, 1 C--H        & $-1$ & anillos externos \\
C=N imina         & 1 C--H, 1 C=N, 1 C--C & $+1$ & enlace \ce{CH=N} \\
CH$_2$ amina      & 2 C--H, 1 C--N, 1 C--C & $-1$ & puente \ce{NH-CH2} \\
\hline
\end{tabular}
\normalsize
\end{center}

\textbf{Cambio por sustitución.} Reemplazar un H aromático por un F 
sube el n.o. de ese carbono de $-1$ a $+1$, un salto de $2$ electrones. 
Es el mismo patrón de la serie \ce{CH4} $\to$ \ce{CO2}, pero aplicado a 
un solo carbono del anillo. En un COF, esto permite modular las 
propiedades electrónicas de un nodo específico sin tocar el resto de la 
red, porque los enlaces C--C (que conectan los nodos entre sí) aportan 
$0$ al n.o. y aíslan el efecto.

\textbf{Suma de carbonos en el fragmento.} Para verificar la coherencia 
del balance, se suman los n.o. de todos los carbonos del fragmento y se 
compara con la carga total, que es $0$ porque la molécula es neutra. El 
resto de los átomos (\ce{H}, \ce{F}, \ce{N}) compensan esa suma:

\begin{center}
\begin{tabular}{lc}
\hline
\textbf{Grupo} & \textbf{Aporte} \\
\hline
Carbonos (suma)   & valor neto \\
Hidrógenos        & $+1$ cada uno \\
Flúor             & $-1$ cada uno \\
Nitrógeno         & $-3$ en imina, $-2$ en amina \\
\hline
Total             & $0$ \\
\hline
\end{tabular}
\end{center}

\textbf{Nota:} el esquema mostrado es un \emph{fragmento modelo}, no un COF 
completo. Un COF real es una red extendida en dos o tres dimensiones donde 
cada nodo se conecta a tres o más vecinos mediante enlaces imina 
(\ce{-CH=N-}), formando una estructura porosa infinita. Aquí se ha cerrado 
el macrociclo para poder mostrarlo como una sola molécula y aplicar el 
análisis de n.o. Los F en los bordes representan sustituyentes que en el 
material real estarían apuntando hacia otros nodos o hacia el espacio 
de poro.

\phantomsection
\addcontentsline{toc}{subsubsection}{Clase 11 -- 01/10/2026: Nivel de oxidación global}
\subsubsection*{\textcolor{blueunam}{Clase 11 -- 01/10/2026: Nivel de oxidación global}}

Además del n.o. de un carbono individual, se puede asignar un 
\textbf{nivel de oxidación global} a una molécula completa. Este número 
caracteriza el estado redox de la molécula entera y sirve para rastrear 
cambios globales en una transformación, no de un sitio específico.

\textbf{Heteroátomo:} átomo distinto de C e H dentro de una molécula 
orgánica. Los más comunes son \ce{O}, \ce{N}, \ce{S} y los halógenos 
(\ce{F}, \ce{Cl}, \ce{Br}, \ce{I}). Introducen diferencia de 
electronegatividad en el esqueleto carbonado, y cada uno cuenta como $+1$ 
en la fórmula.

\textbf{Fórmula general:}
\begin{equation}
\text{n.o. global} = \text{heteroátomos} + \text{anillos} + \text{enlaces } \pi
\end{equation}

\textbf{Cómo se cuenta cada término:}
\begin{itemize}
    \item \textbf{Heteroátomos:} número total de heteroátomos presentes.
    \item \textbf{Anillos:} número de anillos en la molécula.
    \item \textbf{Enlaces $\pi$:} número de dobles enlaces. Los aromáticos cuentan sus dobles individuales.
\end{itemize}

\textbf{Interpretación:}
\begin{itemize}
    \item Un n.o. global \textbf{más bajo} indica una molécula más reducida.
    \item Un n.o. global \textbf{más alto} indica una molécula más oxidada.
    \item Si en una transformación el n.o. global \textbf{no cambia}, el proceso 
    es intramolecular (reorganización interna) y no hay redox neto.
    \item Si el n.o. global \textbf{cambia}, hay redox. La magnitud del cambio 
    indica cuántos electrones se transfieren y si la conversión puede darse 
    en un solo paso o requiere intermediarios (radicales, carbocationes, etc.).
\end{itemize}

\textbf{Ejemplo.} La molécula \chemfig{*6(=-=(-=[:-45])-=-)} (ciclohexadieno 
con sustituyente vinilo) tiene:
\begin{itemize}
    \item Heteroátomos: $0$
    \item Anillos: $1$
    \item Enlaces $\pi$: $4$ (tres del anillo aromático y uno del vinilo)
\end{itemize}
\[
\text{n.o. global} = 0 + 1 + 4 = 5
\]

Al comparar con otra molécula o con el producto de una transformación, la 
diferencia entre los dos n.o. globales da directamente los electrones 
transferidos en el proceso global.

\textbf{Ejemplo.} Pinacol $\to$ pinacolona. \\

\begin{center}
\schemestart
\chemname{\chemfig{CH_3-C^{\color{catred}{+}}(-[:90]OH)(-[:270]CH_3)-C^{\color{catred}{+}}(-[:90]CH_3)(-[:0]OH)(-[:270]CH_3)}}{2,3-dimetil-2,3-butanodiol}
\arrow{->}
\chemname{\chemfig{CH_3-C^0(-[:90]CH_3)(-[:270]CH_3)-C^{\color{catred}{2+}}(-[:90]CH_3)(=[:270]O)}}{3,3-dimetil-2-butanona}
\schemestop
\end{center}

\begin{center}
\renewcommand{\arraystretch}{1.4}
\begin{tabular}{lcccc}
\hline
\textbf{Molécula} & \textbf{Heteroátomos} & \textbf{Anillos} & \textbf{Enlaces }$\boldsymbol{\pi}$ & \textbf{n.o. global} \\
\hline
Pinacol    & $2$ (\ce{O}) & $0$ & $0$ & $2$ \\
Pinacolona & $1$ (\ce{O}) & $0$ & $1$ (\ce{C=O}) & $2$ \\
\hline
\end{tabular}
\end{center}

El n.o. global es igual en reactivo y producto, así que la transformación 
no es redox neta: es una reorganización intramolecular.

\textbf{Ejemplo.} Fotoreducción de benzofenona con 2-propanol. \\

Benzofenona + 2-propanol $\longrightarrow$ Benzopinacol + Propanona \\[8pt]

\begin{center}
\begingroup
\setchemfig{atom sep=2em}
\schemestart
\chemname{\chemfig{O=[:270]C^{\color{catred}{+}}(-[:225]*6(=-=-=-))(-[:-45]*6(=-=-=-))}}{Benzofenona}
\+
\chemname{\chemfig{C^0(-[:90]OH)(-[:-45]CH_3)(-[:225]CH_3)}}{2-propanol}
\arrow{->}
\chemname{\chemfig{C^{\color{catred}{+}}(-[:135]*6(=-=-=-))(-[:225]*6(=-=-=-))(-[:270]OH)-C^{\color{catred}{+}}(-[:45]*6(=-=-=-))(-[:-45]*6(=-=-=-))(-[:90]OH)}}{Benzopinacol}
\+
\chemname{\chemfig{C^{\color{catred}{2+}}(=[:90]O)(-[:-45]CH_3)(-[:225]CH_3)}}{Propanona}
\schemestop
\endgroup
\end{center}

\textbf{Nombres IUPAC:}
\begin{itemize}
    \item Benzofenona: difenilmetanona
    \item 2-propanol: propan-2-ol
    \item Benzopinacol: 1,1,2,2-tetrafeniletano-1,2-diol
    \item Propanona: propan-2-ona (acetona)
\end{itemize}

\textbf{Balanceo:}
\[
\ce{2 (C6H5)2C=O + (CH3)2CHOH -> (C6H5)2C(OH)-C(OH)(C6H5)2 + (CH3)2C=O}
\]

\textbf{Suma de n.o. global:}
\begin{center}
\renewcommand{\arraystretch}{1.4}
\begin{tabular}{lcccc}
\hline
\textbf{Molécula} & \textbf{Heteroátomos} & \textbf{Anillos} & \textbf{Enlaces }$\boldsymbol{\pi}$ & \textbf{n.o. global} \\
\hline
Benzofenona ($\times 2$) & $2$ & $4$ & $14$ & $20$ \\
2-propanol & $1$ & $0$ & $0$ & $1$ \\
\hline
\textbf{Total reactivos} & & & & $21$ \\
\hline
Benzopinacol & $2$ & $4$ & $12$ & $18$ \\
Propanona & $1$ & $0$ & $1$ & $2$ \\
\hline
\textbf{Total productos} & & & & $20$ \\
\hline
\end{tabular}
\end{center}

Las sumas no coinciden ($21 \neq 20$), así que el proceso sí es redox 
neto. La benzofenona se reduce (baja de $20$ a $18$ al convertirse en 
benzopinacol) y el 2-propanol se oxida (sube de $1$ a $2$ al convertirse 
en propanona). La diferencia es de \textbf{1 electrón}, no de $2$. Un 
cambio impar no es compatible con una transferencia en pares como en los 
procesos iónicos, así que la reacción debe pasar por un intermediario 
radicalario: la transferencia ocurre de uno en uno.

\phantomsection
\addcontentsline{toc}{subsubsection}{Balanceo de reacciones orgánicas (Método del [H])}
\subsubsection*{\textcolor{blueunam}{Balanceo de reacciones orgánicas (Método del [H])}}

En química orgánica, el balanceo de reacciones redox no suele hacerse contando electrones libres, sino utilizando el concepto de \textbf{hidrógeno naciente} o \textbf{[H]}. Este método es una adaptación directa del método del ion-electrón, pero resulta mucho más intuitivo para moléculas con esqueletos carbonados.

\textbf{El equivalente químico:} un \ce{[H]} representa formalmente un protón y un electrón.
\[
\ce{[H] \equiv H+ + e-}
\]
Por lo tanto, transferir dos \ce{[H]} es exactamente lo mismo que transferir dos electrones.

\textbf{Especies espectadoras y herramientas de balanceo.}
Antes de aplicar el método, es crucial identificar a los \textbf{iones espectadores}. En las reacciones redox, iones como el potasio (\ce{K+}) o el sodio (\ce{Na+}) no participan en la transferencia de electrones; su número de oxidación no cambia. Por eso se ignoran en las semirreacciones y se trabajan solo con las especies que sí se oxidan o reducen (como el dicromato, \ce{Cr2O7^{2-}}). Al final, los espectadores vuelven a aparecer para formar las sales.

\textbf{Nota importante:} a diferencia de los iones espectadores, los átomos de hidrógeno y oxígeno \textbf{nunca} se ignoran. El agua (\ce{H2O}), los protones (\ce{H+}) y los hidrógenos nacientes (\ce{[H]}) son las herramientas con las que balanceamos la masa y las cargas.

\textbf{Reglas generales para el balanceo con [H]:}
\begin{enumerate}
    \item \textbf{Identificar la oxidación y la reducción:} 
    \begin{itemize}
        \item Oxidación: la molécula gana oxígeno, halógenos, o pierde hidrógenos.
        \item Reducción: la molécula pierde oxígeno, halógenos, o gana hidrógenos.
    \end{itemize}
    
    \item \textbf{Balancear los átomos (esqueleto):} asegurar que el número de carbonos, halógenos, nitrógenos y otros átomos se conserve en ambos lados. \\
    \textit{Tip:} para no arrastrar moléculas grandes, escribe la fórmula molecular o condensada (ej. \ce{C6H5Br} en lugar de dibujar todo el anillo).
    
    \item \textbf{Balancear los oxígenos con \ce{H2O}:} añadir moléculas de agua donde falte oxígeno.
    
    \item \textbf{Balancear los hidrógenos con [H]:} añadir \ce{[H]} en el lado que necesite hidrógenos para igualar la ecuación. Cada \ce{[H]} equivale formalmente a \ce{H+ + e-}.
    
    \item \textbf{Balancear la carga:} en medio ácido, añadir \ce{H+} del lado que tenga la carga más negativa. En medio básico, añadir \ce{OH-} del lado que tenga la carga más positiva.
    
    \item \textbf{Igualar los [H] (electrones):} multiplicar las semirreacciones por los coeficientes necesarios para que los \ce{[H]} ganados en la reducción sean iguales a los perdidos en la oxidación.
    
    \item \textbf{Sumar y cancelar:} sumar las semirreacciones, cancelar los términos semejantes (\ce{[H]}, \ce{H2O}, \ce{H+}) y obtener la ecuación iónica neta.
    
    \item \textbf{Coherencia química y ajuste final:} agrupar los iones para formar las sales moleculares. Luego, verificar que los iones espectadores (como \ce{K+}, \ce{Na+}) y sus contraiones (como \ce{SO4^{2-}}) coincidan en número. Si no coinciden, ajustar los coeficientes de las sales. Finalmente, comprobar que todos los átomos (C, H, O, metales, etc.) estén balanceados.
\end{enumerate}

\vspace{0.5cm}

\textbf{Ejemplo: Oxidación de etanol a ácido acético con dicromato de potasio en medio ácido.}

\begin{center}
\begingroup
\small
\setchemfig{atom sep=1.6em}
\schemestart
\chemfig{CH_3-CH_2-OH}
\+
\ce{K2Cr2O7}
\+
\ce{H2SO4}
\arrow{->}
\chemfig{CH_3-COOH}
\+
\ce{Cr2(SO4)3}
\+
\ce{K2SO4}
\+
\ce{H2O}
\schemestop
\endgroup
\end{center}

\textbf{Paso 1: Semirreacción de Oxidación.}
Usamos la fórmula general \ce{C2H6O} para el etanol y \ce{C2H4O2} para el ácido acético. 

Primero, el esqueleto básico balanceado (carbonos):
\[
\ce{C2H6O -> C2H4O2}
\]

Segundo, balanceamos los oxígenos añadiendo agua en los reactivos (extremo izquierdo):
\[
\ce{H2O + C2H6O -> C2H4O2}
\]

Tercero, balanceamos los hidrógenos añadiendo \ce{[H]} en los productos (extremo derecho):
\begin{equation*}
\ce{H2O + C2H6O -> C2H4O2 + 4[H]} \tag*{\textit{(Oxidación)}}
\end{equation*}
\textit{Nota: el etanol pierde 4 \ce{[H]} (4 electrones).}

\vspace{0.5cm}

\textbf{Paso 2: Semirreacción de Reducción.}
El dicromato (\ce{Cr2O7^{2-}}) se reduce a cromo (III). 

Primero, el esqueleto básico balanceado (cromos):
\[
\ce{2Cr2O7^{2-} -> 4Cr^{3+}}
\]

Segundo, balanceamos los oxígenos añadiendo agua (extremo derecho):
\[
\ce{2Cr2O7^{2-} -> 4Cr^{3+} + 14H2O}
\]

Tercero, balanceamos la carga. Del lado de los reactivos tenemos una carga de $-4$ y del lado de los productos una carga de $+12$. Para igualar en medio ácido, añadimos $16H^+$ al extremo izquierdo de los reactivos:
\[
\ce{16H+ + 2Cr2O7^{2-} -> 4Cr^{3+} + 14H2O}
\]

Cuarto, verificamos los hidrógenos. En los reactivos hay $16$ hidrógenos y en los productos hay $28$ (por las $14H_2O$). Para balancear los $12$ hidrógenos faltantes, añadimos $12[H]$ al extremo izquierdo:
\begin{equation*}
\ce{12[H] + 16H+ + 2Cr2O7^{2-} -> 4Cr^{3+} + 14H2O} \tag*{\textit{(Reducción)}}
\end{equation*}
\textit{Nota: el cromo pasa de n.o. +6 a +3, ganando 12 electrones en total.}

\vspace{0.5cm}

\textbf{Paso 3: Igualar los [H], cancelar e integrar.}
La oxidación produce \(4 \ce{[H]}\) y la reducción consume \(12 \ce{[H]}\). Para igualar los electrones, multiplicamos únicamente la semirreacción de oxidación por 3:
\[
\ce{3H2O + 3C2H6O -> 3C2H4O2 + 12[H]}
\]

Sumamos ambas semirreacciones y cancelamos los términos semejantes algebraicamente:
\begin{itemize}
    \item Los \(12 \ce{[H]}\) están en ambos lados, así que se eliminan por completo (quedan a 0).
    \item Los \(\ce{H2O}\) no se eliminan del todo: hay 3 a la izquierda y 14 a la derecha. Al restar \(14 - 3 = 11\), quedan 11 moléculas de agua en el lado de los productos.
\end{itemize}

\[
\ce{\cancelto{0}{12[H]} + 16H+ + \cancelto{0}{3H2O} + 3C2H6O + 2Cr2O7^{2-} -> 3C2H4O2 + 4Cr^{3+} + \cancelto{11}{14H2O} + \cancelto{0}{12[H]}}
\]

Agrupando los reactivos, la ecuación iónica neta es:
\[
\ce{16H+ + 3C2H6O + 2Cr2O7^{2-} -> 3C2H4O2 + 4Cr^{3+} + 11H2O}
\]
\textit{Comprobación de cargas:} a la izquierda hay $+12$ (del \(16H^+\) menos el \(-4\) de los dos dicromatos). A la derecha hay $+12$ (de los cuatro \(Cr^{3+}\)). Las cargas están balanceadas.

\vspace{0.5cm}

\textbf{Paso 4: Coherencia química (formar sales y ajuste final).}
Los $16 \ce{H+}$ de la ecuación iónica provienen del ácido sulfúrico (\ce{H2SO4}). Como es un ácido diprótico (aporta dos \(H^+\) por molécula), se necesitan $8 \ce{H2SO4}$ para obtener los 16 protones. 

Al armar la ecuación molecular juntando los iones, escribimos tentativamente las sales:
\[
\ce{3C2H6O + 2K2Cr2O7 + 8H2SO4 -> 3C2H4O2 + 2Cr2(SO4)3 + K2SO4 + 11H2O}
\]

\textbf{Ajuste final:} si observamos el potasio, del lado de los reactivos tenemos $4 \ce{K}$ (por los $2\ce{K2Cr2O7}$), pero del lado de los productos solo hay $2 \ce{K}$ (por el $\ce{K2SO4}$). Para balancear el potasio, el coeficiente del sulfato de potasio debe ser 2. Este ajuste también balancea los sulfatos (S) y oxígenos (O) de esa sal:
\[
\ce{3C2H6O + 2K2Cr2O7 + 8H2SO4 -> 3C2H4O2 + 2Cr2(SO4)3 + 2K2SO4 + 11H2O}
\]

\textbf{Comprobación final de átomos:}
\begin{itemize}
    \item \textbf{C:} $3 \times 2 = 6$ a la izquierda; $3 \times 2 = 6$ a la derecha.
    \item \textbf{H:} $3(6) + 8(2) = 34$ a la izquierda; $3(4) + 11(2) = 34$ a la derecha.
    \item \textbf{O:} $3(1) + 2(7) + 8(4) = 49$ a la izquierda; $3(2) + 2(12) + 2(4) + 11(1) = 49$ a la derecha.
    \item \textbf{K:} $2 \times 2 = 4$ a la izquierda; $2 \times 2 = 4$ a la derecha.
    \item \textbf{Cr:} $2 \times 2 = 4$ a la izquierda; $2 \times 2 = 4$ a la derecha.
    \item \textbf{S:} $8 \times 1 = 8$ a la izquierda; $2(3) + 2(1) = 8$ a la derecha.
\end{itemize}
La ecuación está totalmente balanceada.

\vspace{0.5cm}

\textbf{Información que aporta el balanceo:}
\begin{itemize}
    \item \textbf{Electrones intercambiados:} Se transfirieron 12 electrones en total (evidenciados por los 12 \ce{[H]} que se cancelaron en el Paso 3).
    \item \textbf{¿Quién se oxida?} El etanol (\ce{C2H6O}), porque pierde hidrógenos y gana oxígeno.
    \item \textbf{¿Quién se reduce?} El dicromato (\ce{Cr2O7^{2-}}), porque el cromo pasa de n.o. +6 a +3.
    \item \textbf{Agente reductor:} El etanol (es la especie que se oxida y provoca que la otra se reduzca).
    \item \textbf{Agente oxidante:} El dicromato (es la especie que se reduce y provoca que la otra se oxide).
    \item \textbf{Ecuación molecular balanceada:} 
    \[
    \ce{3C2H6O + 2K2Cr2O7 + 8H2SO4 -> 3C2H4O2 + 2Cr2(SO4)3 + 2K2SO4 + 11H2O}
    \]
\end{itemize}